% TOP_PROBLEM: {"id": 1108, "title": "Andrews-Curtis Conjecture", "category_id": 4, "category": {"id": 4, "name": "algebra", "display_name": "Algebra", "description": "Group theory, ring theory, field theory, and algebraic structures.", "slug": "algebra", "order_index": 4, "created_at": "2026-07-31T15:26:25.671Z"}, "status": "open"}
% FIRST_POSTED: 2026-09-27
% !TeX program = lualatex
% ATTEMPT_STATUS: unresolved
% Model: GPT-6 Astra Ultra; continuation dated 27 September 2026.
\documentclass[11pt]{article}
\usepackage[margin=25mm]{geometry}
\usepackage{fontspec}
\setmainfont{Latin Modern Roman}
\usepackage{amsmath,amssymb,mathtools}
\usepackage{unicode-math}
\setmathfont{Latin Modern Math}
\usepackage{xurl,hyperref}
\hypersetup{colorlinks=true,urlcolor=blue}
\setcounter{secnumdepth}{0}
\setlength{\parindent}{0pt}
\setlength{\parskip}{5pt}
\setlength{\emergencystretch}{3em}
\begin{document}
\section*{\texorpdfstring{Andrews-Curtis Conjecture}{Andrews-Curtis Conjecture}}\label{1108}
\subsection{Definitions and mathematical statement}
A balanced finite presentation is
\[
 \langle x_1,\ldots,x_n\mid r_1,\ldots,r_n\rangle,
\]
with equally many generators and relators, the latter words in the free group $F_n$. It presents the trivial group when their normal closure is all of $F_n$. The allowed operations on the relator tuple are inversion $r_i\mapsto r_i^{-1}$, multiplication $r_i\mapsto r_ir_j$ for $i\ne j$, interchange of relators, and conjugation $r_i\mapsto w r_iw^{-1}$ by a word $w\in F_n$.

The conjecture asks whether every such trivial presentation can be transformed in finitely many allowed moves to $(x_1,\ldots,x_n)$. The generator set stays fixed. Adding auxiliary generator--relator pairs is a stabilization and is not an allowed extra operation in this unstabilized problem.
\par\smallskip\textit{Formulation and concept references:} \href{https://link.springer.com/article/10.1007/s10817-026-09759-8}{[S1]}.

\subsection{Short English statement}
Can every balanced presentation of the trivial group be simplified to the standard presentation using only the prescribed relator moves?
\subsection{Research attempt}

\paragraph{Scope, source audit and certificate standard.}
The generators remain fixed throughout this attempt. Every asserted
simplification is a composition of relator inversion, multiplication
by a different relator, interchange, and conjugation by a word in the
fixed free group. Introducing an auxiliary generator with a matching
relator would address the stabilized problem, which is not the target.

The supplied Fairbank--Lisitsa--Vernitski article [S1], published
8 July 2026, studies the conjecture through classification and guided
search and states it as unresolved. Its computational methods motivate
checking explicit move certificates, but a classifier output or a failed
search is not a group-theoretic proof. Exhaustively enumerating finite
move sequences is a semidecision procedure for finding a successful
simplification; no decision procedure for the negative case follows
from that enumeration alone.

The attempt first separates abelian obstructions from the free-group
problem, proves a relator-erasure lemma, and uses it to solve explicit
triangular families and the rank-two case with a primitive relator.
It then constructs exact certificates and identifies the remaining
normal-generation gap. These are elementary standard mechanisms, not
a claim of new progress on the general conjecture.

\paragraph{Normal closure is an invariant of every allowed move.}
For a relator tuple $r=(r_1,\ldots,r_n)$ put
$N(r)=\langle\!\langle r_1,\ldots,r_n\rangle\!\rangle\triangleleft F_n$.
Inverting or conjugating one relator leaves this normal subgroup
unchanged, as does interchange. For multiplication, the new relator
$r_ir_j$ lies in the old normal closure, and conversely
$r_i=(r_ir_j)r_j^{-1}$ lies in the new one. Thus this move also preserves
$N(r)$. Every allowed operation is reversible by finitely many allowed
operations; for example right multiplication by $r_j^{-1}$ is achieved
by inverting $r_j$, multiplying, and restoring $r_j$.

It follows that a certificate ending at $(x_1,\ldots,x_n)$ proves both
the triviality of the presented group and its AC equivalence to the
standard tuple. A proof of triviality by some other means need not
produce such a certificate: it only identifies the normal subgroup,
not the connected component in the graph of allowed tuples.

\paragraph{The exact abelian obstruction and its limitation.}
For a word $w\in F_n$, let $e_j(w)$ be its exponent sum in $x_j$.
The matrix
\[
 E(r)=(e_j(r_i))_{1\leq i,j\leq n}
\]
records the relations in the abelianization. If the presented group is
trivial, its abelianization is zero, so the rows span $\mathbb Z^n$.
For a square integral matrix this is equivalent to $\det E(r)=\pm1$.
Inversion negates a row, relator multiplication adds a row, interchange
swaps rows, and conjugation leaves every row unchanged.

\textbf{Lemma 456.1.} If $\det E(r)=\pm1$, allowed relator moves
transform $r$ to a tuple with exponent matrix $I_n$.

\emph{Proof.}
Integer elementary row operations reduce a unimodular matrix to the
identity by the Euclidean algorithm. Row negation and interchange are
allowed moves. Adding a positive integer multiple of row $j$ to row
$i$ is repeated multiplication of $r_i$ by $r_j$. For a negative
multiple, first invert $r_j$, perform the positive number of
multiplications, and restore it. These are actual free-group operations;
their exponent matrices are the desired row operations.\hfill$\square$

After this normalization, each relator has the form
$r_i=x_ic_i$ with $c_i\in[F_n,F_n]$. The lemma does not remove the
$c_i$. Treating exponent-sum cancellation as free reduction would
erase exactly the information that remains difficult.

A concrete balanced presentation shows why unimodularity alone cannot
even certify triviality. Take
\[
 r_1=(ab)^2a^{-3},\qquad r_2=a^3b^{-5}.
\]
Its exponent matrix and determinant are
\[
 \begin{pmatrix}-1&2\\3&-5\end{pmatrix},\qquad \det=-1.
\]
Nevertheless it has a nontrivial permutation representation. With
permutations composed right to left, put
\[
 a=(1\ 2\ 3),\qquad b=(1\ 4\ 3\ 2\ 5).
\]
Then $a^3=b^5=1$ and $ab=(1\ 4)(2\ 5)$, so both relators evaluate
to the identity while $a$ does not. Thus the group is not trivial.
We do not need to identify that group to obtain the obstruction.
By the normal-closure invariant, this tuple cannot simplify to $(a,b)$.
The example is outside the hypothesis of the conjecture and only
refutes the proposed intermediate test ``determinant one implies trivial''.

\paragraph{Erasing occurrences using an already isolated generator.}
\textbf{Lemma 456.2.} Suppose one relator is exactly $x_j$. Keeping
it equal to $x_j$ at the end of each operation block, one can replace
any other relator $w$ by the word obtained by setting $x_j=1$ in $w$,
using only allowed moves.

\emph{Proof.}
Choose one occurrence and write the current word as
$w=u x_j^{\epsilon}v$, where $\epsilon\in\{1,-1\}$.
Right multiplication by $v^{-1}x_j^{-\epsilon}v$ gives
\[
 u x_j^\epsilon v\,v^{-1}x_j^{-\epsilon}v=uv
\]
in the free group. To perform this multiplication using the stored
relator $x_j$, invert that relator if needed, conjugate it by $v^{-1}$,
multiply the target relator by it, undo the conjugation, and undo
the possible inversion. There are at most five permitted moves.
The stored relator is restored to $x_j$; all other relators except
the target are unchanged. Each block removes the chosen occurrence,
and free reduction can only shorten the remaining word. Repetition
terminates and yields exactly the image of $w$ under $x_j\mapsto1$.
\hfill$\square$

The words $u,v$ in this proof are fixed free-group words at the moment
of the operation. Conjugating by $v^{-1}$ is allowed even when $v$
contains other generators. There is no introduction of a generator
representing $v$, no new relator and no deletion of a generator.
This distinction makes the lemma valid for the unstabilized problem.

\paragraph{A recursive sufficient criterion and triangular families.}
Suppose $r_1=x_1$, and let $\overline r_i$ be obtained from $r_i$
by setting $x_1=1$ for $i\geq2$. If
$(\overline r_2,\ldots,\overline r_n)$ can be AC simplified in
$F(x_2,\ldots,x_n)$, Lemma 456.2 first erases all occurrences of
$x_1$ in those relators. The smaller-rank certificate then operates
on the remaining slots, leaving the first slot equal to $x_1$.
Every conjugating word in that certificate is also a word in $F_n$.
Thus the original tuple simplifies without a stabilization.

In particular, consider
\[
 r_i=u_i x_i^{\epsilon_i}v_i\quad(1\leq i\leq n),\qquad
 u_i,v_i\in F(x_{i+1},\ldots,x_n),\quad\epsilon_i\in\{1,-1\}.
 \tag{456.1}
\]
The last relator is $x_n^{\epsilon_n}$ and can be inverted to $x_n$.
Use it to erase $x_n$ from all earlier slots. The penultimate slot
then becomes $x_{n-1}^{\epsilon_{n-1}}$; repeat downward. This is a
finite, constructive proof of AC simplifiability for the entire family
(456.1), not merely a proof that its quotient group is trivial.

An occurrence-count bound makes the procedure explicit. If the initial
total length is $N$, erasing the later generators requires at most
$N$ erasure blocks, since after every block the actual relator tuple
has shorter target words and the stored relator is restored. At most
$n$ final inversions are needed. Thus the displayed implementation
uses at most $5N+n$ moves. Intermediate conjugated helper words may
be long; this estimate is a bound on the number of allowed moves,
where a conjugation by an arbitrary word counts as one move. It is
not a uniform bound for arbitrary trivial presentations.

\paragraph{Primitive relators and the role of ambient automorphisms.}
A word is primitive if it belongs to a free basis. This is much
stronger than its exponent vector having relatively prime entries.
For example, $a^2b^3$ has exponent vector $(2,3)$ but is not primitive
in $F(a,b)$. If it were primitive, quotienting by its normal closure
would give a free cyclic group. Instead the assignments
$a\mapsto(1\ 2)$ and $b\mapsto(1\ 2\ 3)$ satisfy the relation
and generate a nonabelian group. Hence that quotient is not cyclic.

We now justify a familiar change-of-basis argument without silently
adding generator changes to the allowed move list. Let
$\phi\in\operatorname{Aut}(F_n)$. Applying $\phi$ to every word
in a move certificate produces another allowed certificate: inversion
and multiplication commute with $\phi$, and conjugation by $w$
becomes conjugation by $\phi(w)$.

We also use Nielsen's classical theorem that any free basis is connected
to the standard one by inversions, interchanges and elementary
multiplications. These operations on a basis tuple are allowed
relator moves. If a chosen convention uses left multiplication
$r_i\mapsto r_jr_i$, it is simulated by conjugating $r_i$ by $r_j$
and then right multiplying by $r_j$:
$(r_jr_ir_j^{-1})r_j=r_jr_i$.
Thus every free basis, viewed as a relator tuple, is AC equivalent
to the standard tuple.

It follows that a tuple is AC trivial if and only if its image under
$\phi$ is AC trivial. For the reverse direction, apply $\phi^{-1}$
to a certificate ending at the standard tuple; it now ends at the
free basis $(\phi^{-1}(x_i))$, which has its own Nielsen certificate
to the standard tuple. This proves invariance of the property, while
the actual move sequences still keep the ambient generators fixed.

\textbf{Proposition 456.3.} A balanced trivial presentation on two
generators with at least one primitive relator is AC trivial.

\emph{Proof.}
After interchange and the justified automorphism reduction, assume
the first relator is $a$. Erase $a$ from the second relator using
Lemma 456.2. It becomes $b^m$ for an integer $m$.
The quotient group is unchanged and is now
$\langle b\mid b^m\rangle$. Triviality forces $m=1$ or $m=-1$:
$m=0$ gives an infinite cyclic group and $|m|>1$ a nontrivial finite
cyclic group. Invert the second relator if necessary to obtain $(a,b)$,
and translate the certificate back as above.\hfill$\square$

For $n>2$, the analogous argument reduces a primitive-relator case
to the conjecture in rank $n-1$; it does not unconditionally settle
all those higher-rank cases. In rank one the conjecture itself is
immediate: a single relator in $F_1$ is $x^m$, and its quotient is
trivial exactly for $m=\pm1$.

\paragraph{A complete small certificate.}
The tuple $(ab,bab)$ presents the trivial group and admits the
following unstabilized certificate. Each line is a tuple of freely
reduced words, and the indicated move changes just one slot:
\[
\begin{array}{c|l}
 (ab,bab)&\text{initial tuple}\\
 (b^{-1}a^{-1},bab)&\text{invert slot 1}\\
 (b^{-1}a^{-1},b)&\text{slot 2 times slot 1}\\
 (ab,b)&\text{invert slot 1}\\
 (ab,b^{-1})&\text{invert slot 2}\\
 (a,b^{-1})&\text{slot 1 times slot 2}\\
 (a,b)&\text{invert slot 2}.
\end{array}
\]
The same mechanism handles $(ab^m,bab^m)$ for every integer $m$:
the second relator times the inverse of the first is $b$, after which
the powers of $b$ in the first can be erased. The short example is
included to make multiplication orientation and helper restoration
auditable; it is not put forward as a difficult former candidate.

\paragraph{What a search can and cannot prove.}
Represent words by signed generator indices and perform free reduction
with a stack: append a letter unless it cancels the final letter,
in which case remove that final letter. The uniqueness of freely
reduced representatives in a free group makes equality checks exact.
The accompanying verifier checks every transition in the displayed
certificate, tests the erasure block on explicit words, and verifies
generated scramble-and-reverse certificates without changing rank.
It also checks the two permutation diagnostics above.

For an exhaustive search it is enough to permit conjugation by one
generator or its inverse at a time: conjugation by an arbitrary word
is a finite composition of such conjugations, taken in the appropriate
reverse order of its letters. With inversion, multiplication and
interchange this produces a finitely branching move graph. Breadth-first
search eventually finds any finite certificate if one exists.

Failure up to a depth bound only excludes certificates within that
bound. Imposing a total-word-length cutoff additionally removes paths
that must leave the permitted region before returning. A claimed
counterexample would therefore require a proved invariant separating
components, or another global obstruction, not an empty search frontier
inside a prescribed finite box. Random walks from the standard tuple
produce useful positive test data, but all resulting presentations are
known to be AC trivial by construction; they cannot test the existence
of unknown components containing trivial presentations.

\paragraph{The precise gap after normal generation.}
If $N(r)=F_n$, each generator has an expression
\[
 x_i=\prod_{j=1}^{m_i}w_{ij}r_{k_{ij}}^{\epsilon_{ij}}w_{ij}^{-1}.
\]
This is exactly what membership in a normal closure means. It does
not provide an allowed move replacing one relator by that product
while retaining a usable copy of every relator appearing in the
expression. Such a replacement is not one of the permitted operations;
attempting to compute all the products can overwrite the slots needed
later. Adding scratch relators or extra generator--relator pairs would
change the problem unless their use is rigorously eliminated.

The erasure lemma resolves this resource issue when one actual generator
is already stored as a relator, and the primitive-relator argument
creates that situation in a controlled special case. No argument here
shows that an arbitrary normally generating balanced tuple can be moved
to one having such a relator. That is the first missing step of the
present strategy. The general unstabilized conjecture remains unresolved.

\subsection{Sources}
\begin{itemize}
\item[S1]Michael Fairbank, Alexei Lisitsa, and Alexei Vernitski, Journal of Automated Reasoning 70 (2026), article 11.
\url{https://link.springer.com/article/10.1007/s10817-026-09759-8}.
\textit{Formulation source.}
\item[E1]Michael Fairbank, Alexei Lisitsa and Alexei Vernitski,
\emph{Probabilistic Automaton Classifier Applied to Examples Related
to the Andrews--Curtis Conjecture}, Journal of Automated Reasoning
70 (2026), article 11, published 8 July 2026; full text of [S1].
\url{https://link.springer.com/article/10.1007/s10817-026-09759-8}.
\textit{Primary Sections 1--2 inspected 27 September 2026. The move
set, stabilization distinction and computational context are retained;
enumerating certificates is described only as a semidecision procedure.}
\end{itemize}
\end{document}
